\subsection{Statements, Basic Blocks, and Execution Edges}

A Python source file is written as a vertical sequence of statements. At first sight, this makes a program look like a simple list: first line, second line, third line, and so on. That view is useful only for the simplest programs. As soon as the program contains a condition, a loop, or an exception, the execution path no longer follows the written text in a single uninterrupted direction.

A more accurate model is a control flow graph. In this model, the program is divided into pieces of code that can execute straight through without an internal jump. Such a piece is usually called a basic block. Between these blocks there are execution edges. An edge means: after this block finishes, execution may continue in that other block.

This does not mean that the beginner must draw a graph for every program. The point is mental, not decorative. A program is not merely text. It is a set of possible routes through text.

\begin{quote}
Source code is linear as written. Execution is graph-shaped as soon as the program contains alternatives, repetition, or failure paths.
\end{quote}

This matters especially for programmers coming from C, because the same deep idea exists there too. In C, \texttt{if}, \texttt{while}, \texttt{for}, \texttt{return}, \texttt{break}, \texttt{continue}, and \texttt{goto} all alter the possible route of execution. Python has different surface syntax, but the control problem is the same: the runtime must know which statement may execute next.

The following very small program is written as a sequence:

\begin{verbatim}
name = "George"
score = 42

print(f"name: {name}")
print(f"score: {score}")
print("These pretzels are making me thirsty.")
\end{verbatim}

Possible output:

\begin{verbatim}
name: George
score: 42
These pretzels are making me thirsty.
\end{verbatim}

Here, the control flow is still simple. Each statement falls into the next one. There is only one ordinary route:

\begin{verbatim}
statement 1 -> statement 2 -> statement 3 -> statement 4 -> end
\end{verbatim}

Once branching, looping, and exceptions appear, that simple route becomes only one special case of the larger execution model.

\subsubsection{Sequential Flow: The Default Fall-Through Path from One Statement to the Next}

Sequential flow is the default case. If nothing interrupts execution, Python executes one statement and then continues with the next statement in the same block.

This ordinary movement from one statement to the next is often called fall-through. The name is useful because it describes the simplest edge in the control flow graph: execution falls from the current statement into the next statement below it.

\begin{verbatim}
quote = "Nobody expects the Spanish Inquisition!"
count = 3
result = count + 39

print(f"quote: {quote}")
print(f"count: {count}")
print(f"result: {result}")
\end{verbatim}

Possible output:

\begin{verbatim}
quote: Nobody expects the Spanish Inquisition!
count: 3
result: 42
\end{verbatim}

The route is direct:

\begin{verbatim}
quote assignment
    -> count assignment
    -> result assignment
    -> first print
    -> second print
    -> third print
\end{verbatim}

There is no condition to test, no loop target to revisit, and no handler to enter. Every statement is reached exactly once, in written order.

This is the simplest kind of basic block: a straight-line segment of code.

\begin{quote}
A basic block is a piece of code whose internal statements execute in sequence, without an internal branch target or jump out of the middle.
\end{quote}

In real Python programs, a basic block may be very short. Sometimes it is only one statement. This is normal. The important property is not size. The important property is that once execution enters the block, ordinary execution proceeds through it in order until the block ends or an interrupting control transfer occurs.

Consider this block inside a future \texttt{if} statement:

\begin{verbatim}
print("checking soup status")
soup_available = False
print(f"soup_available: {soup_available}")
\end{verbatim}

Those three statements form a straight route. If execution enters the block, the ordinary path runs through them from top to bottom.

Sequential flow is therefore the foundation. Branches, loops, and exceptions do not replace sequential flow. They connect sequential blocks in more complex ways.

\subsubsection{Conditional Edges: Branching Execution Paths Created by \texttt{if}, \texttt{elif}, and \texttt{else}}

A conditional statement creates multiple possible outgoing edges from a decision point. Python evaluates a condition, then chooses one continuation path.

In ordinary written form, the code is still vertical:

\begin{verbatim}
score = 42

if score > 90:
    print("excellent")
elif score >= 50:
    print("acceptable")
else:
    print("needs work")

print("grading finished")
\end{verbatim}

Possible output:

\begin{verbatim}
needs work
grading finished
\end{verbatim}

The written order is not the same thing as the execution route. Python does not execute every branch body. It tests the first condition. If that condition is false, it tests the next condition. If all previous tested conditions fail, it enters the \texttt{else} body when one exists.

The control shape is closer to this:

\begin{verbatim}
start
  -> test score > 90
       true  -> print "excellent"   -> after if
       false -> test score >= 50
                    true  -> print "acceptable" -> after if
                    false -> print "needs work" -> after if
  -> print "grading finished"
\end{verbatim}

The \texttt{if}, \texttt{elif}, and \texttt{else} keywords therefore do not merely group text. They define possible execution edges.

A useful way to read the statement is:

\begin{quote}
At this point, execution must choose one of several possible routes. After the chosen route finishes, execution joins again after the conditional structure.
\end{quote}

Only one branch body is selected in a normal \texttt{if}/\texttt{elif}/\texttt{else} chain. This is different from writing separate \texttt{if} statements.

\begin{verbatim}
value = 42

if value > 0:
    print("positive")

if value % 2 == 0:
    print("even")

if value == 42:
    print("the answer")
\end{verbatim}

Possible output:

\begin{verbatim}
positive
even
the answer
\end{verbatim}

Here, there are three separate decision points. All three tests are evaluated, because each \texttt{if} statement starts a new conditional structure.

By contrast:

\begin{verbatim}
value = 42

if value < 0:
    print("negative")
elif value % 2 == 0:
    print("even")
elif value == 42:
    print("the answer")
else:
    print("something else")
\end{verbatim}

Possible output:

\begin{verbatim}
even
\end{verbatim}

The third condition is not reached, even though \texttt{value == 42} is true. The earlier \texttt{elif} branch already captured the route.

This distinction is easier to understand with the graph model. Separate \texttt{if} statements create separate decision nodes. An \texttt{if}/\texttt{elif}/\texttt{else} chain creates one connected decision structure where the first successful branch consumes the path.

\subsubsection{Loop Back-Edges: Repeated Execution Paths Created by \texttt{while} and \texttt{for}}

A loop creates a route that can return to an earlier decision point. This returning route is the essential difference between a loop and a simple conditional branch.

A conditional statement chooses a path and then usually moves forward. A loop may choose a path, execute a block, and then jump back to test or retrieve the next value again.

A \texttt{while} loop repeats as long as its condition remains true:

\begin{verbatim}
n = 1

while n <= 3:
    print(f"n: {n}")
    n = n + 1

print("loop finished")
\end{verbatim}

Possible output:

\begin{verbatim}
n: 1
n: 2
n: 3
loop finished
\end{verbatim}

The graph shape is not a straight line:

\begin{verbatim}
start
  -> set n = 1
  -> test n <= 3
       true  -> print n
             -> increment n
             -> go back to test n <= 3
       false -> print "loop finished"
\end{verbatim}

The edge from the end of the loop body back to the condition is the loop back-edge. Without this edge, there is no repetition.

A \texttt{for} loop has a similar graph shape, but the repeated decision is not written as an explicit boolean condition. Instead, Python asks an iterable object for successive values. Later sections will examine that mechanism through \texttt{iter()}, \texttt{next()}, and \texttt{StopIteration}. For now, the control-flow idea is enough.

\begin{verbatim}
names = ["Jerry", "Elaine", "Kramer"]

for name in names:
    print(f"current name: {name}")

print("all names processed")
\end{verbatim}

Possible output:

\begin{verbatim}
current name: Jerry
current name: Elaine
current name: Kramer
all names processed
\end{verbatim}

Its control shape can be read like this:

\begin{verbatim}
start
  -> get traversal source
  -> ask for next value
       value exists -> bind name
                    -> execute loop body
                    -> go back and ask for next value
       no value     -> continue after loop
\end{verbatim}

This is why loops must be treated as graph structures, not merely repeated text. The loop body is written once, but the execution route may pass through that body many times.

The difference between \texttt{while} and \texttt{for} is therefore not that one loops and the other does not. Both loop. The difference is where the continuation decision comes from.

\begin{itemize}
    \item A \texttt{while} loop repeats by rechecking a condition written by the programmer.
    \item A \texttt{for} loop repeats by asking a traversal object for the next value.
\end{itemize}

Both forms create a back-edge in the control flow graph.

\subsubsection{Exceptional Edges: Non-Local Control Transfers Created by Exceptions}

An exception creates a different kind of execution edge. It does not merely choose between ordinary nearby branches, and it does not simply loop back to an earlier point. It can leave the current expression, the current statement, the current block, and sometimes the current function call, searching for a handler.

This is why exceptions are called non-local control transfers. The next executed statement may not be the next statement in the same block.

Consider this simple program:

\begin{verbatim}
text = "forty-two"

print("before conversion")

number = int(text)

print(f"number: {number}")
print("after conversion")
\end{verbatim}

Possible output:

\begin{verbatim}
before conversion
Traceback (most recent call last):
  ...
ValueError: invalid literal for int() with base 10: 'forty-two'
\end{verbatim}

The assignment to \texttt{number} does not complete. The following two \texttt{print()} calls are not executed. The ordinary fall-through edge is abandoned because \texttt{int(text)} raises a \texttt{ValueError}.

With a handler, the exceptional route is caught:

\begin{verbatim}
text = "forty-two"

print("before conversion")

try:
    number = int(text)
    print(f"number: {number}")
except ValueError as err:
    print("conversion failed")
    print(f"error type: {type(err)}")
    print(f"error message: {err}")

print("after try statement")
\end{verbatim}

Possible output:

\begin{verbatim}
before conversion
conversion failed
error type: <class 'ValueError'>
error message: invalid literal for int() with base 10: 'forty-two'
after try statement
\end{verbatim}

The important fact is not only that an error occurred. The important control-flow fact is that execution jumped from inside the \texttt{try} block to the matching \texttt{except} block. The rest of the ordinary \texttt{try} block was skipped.

The graph shape is approximately:

\begin{verbatim}
start
  -> print "before conversion"
  -> enter try block
  -> attempt int(text)
       success -> bind number
               -> print number
               -> after try statement
       ValueError -> matching except block
                  -> print failure information
                  -> after try statement
\end{verbatim}

This is very different from the C habit of checking returned error codes after each operation. In C, a common pattern is to keep execution local and manually inspect the result:

\begin{verbatim}
/* C-style idea, not Python code */
result = parse_number(text);
if (result.failed) {
    handle_error(result);
}
\end{verbatim}

Python can also use explicit checks, but its exception mechanism gives operations the ability to redirect execution when normal completion is impossible.

The object captured by \texttt{except ValueError as err} is itself a runtime object. It has a type, a textual representation, and methods inherited from the exception hierarchy.

\begin{verbatim}
text = "forty-two"

try:
    number = int(text)
except ValueError as err:
    print(f"type(err): {type(err)}")
    print(f"str(err): {str(err)}")
    print(f"args: {err.args}")
    print(f"has __traceback__: {'__traceback__' in dir(err)}")
\end{verbatim}

Possible output:

\begin{verbatim}
type(err): <class 'ValueError'>
str(err): invalid literal for int() with base 10: 'forty-two'
args: ("invalid literal for int() with base 10: 'forty-two'",)
has __traceback__: True
\end{verbatim}

The \texttt{args} attribute contains the construction arguments stored by the exception object. The \texttt{\_\_traceback\_\_} attribute connects the exception object to traceback information, which records where the failure passed through the active execution frames. Later sections will study this mechanism more carefully.

For now, the essential rule is:

\begin{quote}
An exception adds a hidden possible edge from the place where failure occurs to the nearest matching handler, or farther upward if no local handler exists.
\end{quote}

This completes the first control-flow map:

\begin{itemize}
    \item Sequential statements create fall-through edges.
    \item Conditional statements create branch edges.
    \item Loops create back-edges.
    \item Exceptions create non-local exceptional edges.
\end{itemize}

Once these four edge types are understood, Python source code can be read as more than text. It becomes a set of possible execution routes through runtime state.